\chapter{Discussion} \label{ch:discussion}

This chapter addresses limitations of the methodology, discusses mitigations for
the vulnerabilities identified, and identifies open problems for future work.

% ======================================================================
\section{Limitations} \label{se:limitations}

\subsection{Regtest Environment}

The regtest limitations detailed in Section~\ref{se:regtest_validity} bound the
external validity of all measurements. The most significant gap is the absence of
temporal dynamics: recovery races between the attacker's withdrawal and the
watchtower's response are instantaneous on regtest but uncertain on mainnet, where
block confirmation times are stochastic and mempool congestion affects transaction
propagation. The No Pareto-Optimal Design theorem
(Proposition~\ref{thm:fee_inversion}) is robust to this gap because its
per-threat orderings rest on structural vsize measurements rather than
timing---but the \emph{practical feasibility} of per-UTXO recovery-cost scaling
at a given fee rate depends on how quickly the attacker can fill
blocks with split transactions, which regtest cannot test.

\subsection{Reference Implementations} \label{se:impl_choices}

This work studies covenant \emph{designs} (BIPs), not implementations. However, the
measurements come from specific reference implementations, and implementation choices
affect the measured vsize. The most significant example is OP\_VAULT's fee-input
requirement: the 292\vB{} trigger measurement includes a fee-wallet input that is
inherent to BIP-345's value preservation rule, not an implementation artifact. But
CTV's 3{,}300-sat anchor output and CAT+CSFS's unconstrained recovery leaf are
implementation choices, not BIP requirements. Table~\ref{tab:impl_choices} lists
the design choices in the two author-built implementations (CAT+CSFS and Simplicity)
that could affect measured properties.

\begin{table}[t]
\centering
\caption{Implementation choices in author-built vaults. These are design decisions,
not BIP requirements; alternative designs would produce different measurements.}
\label{tab:impl_choices}
\small
\begin{tabular}{@{}llll@{}}
\toprule
Choice & Covenant & BIP-required? & Alternative \\
\midrule
Fixed destination        & CAT+CSFS & No  & Poelstra's dynamic destination~\cite{Poe21} \\
Unconstrained recovery   & CAT+CSFS & No  & Output-constrained (like Simplicity) \\
MAX\_FEE cap (in-value) & Simplicity & No & Fully dynamic signer-side fee selection without cap \\
3{,}300-sat anchor       & CTV (upstream) & No & Different anchor size \\
\bottomrule
\end{tabular}
\end{table}

\subsection{Rational Adversary Assumption}

The economic analysis assumes a rational adversary who mounts an attack only when
$\text{attack\_cost} < \text{vault\_value}$. This excludes irrational adversaries
(state actors, ideologically motivated attackers) who may attack regardless of cost.
It also excludes sophisticated attackers who combine covenant-level attacks with
network-level attacks (eclipse, time-dilation) to improve their success probability.

\subsection{Threats Not Modeled} \label{se:threats_not_modeled}

The TM1--TM11 threat taxonomy covers covenant-specific attack classes but excludes
several categories that apply equally to all vault designs:

\begin{itemize}
\item \textbf{TM3 race dynamics are network-layer, not Script-layer.}
  Our design-space decomposition
  (Section~\ref{ss:design_space}) does not model the race between
  the attacker's unvault-plus-withdrawal sequence and the
  watchtower's recovery broadcast. Formalising
  $\Pr_D[\text{attacker wins race}]$ would require a separate
  probability space over block-arrival times, mempool propagation
  delays, and detection latency---parameters that depend on the
  deployment environment (mainnet, regtest, signet) rather than on
  BIP script semantics. Existing watchtower and payment-channel
  formalisations in the state-channel literature treat confirmation
  as atomic or assume bounded-delay synchrony; none provides
  off-the-shelf machinery for a pinning-composed race against
  a Bitcoin mempool. We therefore treat TM3 as uniformly
  applicable to all four BIPs at the Script layer. CAT+CSFS's
  dual-binding on the hot leaf is the one design-level mitigation
  in our taxonomy: it reduces the worst-case outcome from theft to
  griefing (see Section~\ref{se:trigger_theft}). A formal network-layer
  race model is left for future work.

\item \textbf{Mining-level attacks.} An attacker with significant hash power could
  withhold blocks to compress CSV timelocks, reducing the watchtower's response
  window. This interacts with TM3 and TM4 but is not covenant-specific.

\item \textbf{Network-level attacks.} Eclipse attacks could prevent a watchtower from
  observing unvault transactions, giving the attacker a head start on the CSV race.

\item \textbf{Watchtower availability.} All vault designs assume the watchtower is
  online during the CSV window. DDoS attacks targeting the watchtower infrastructure
  are a deployment concern orthogonal to covenant design.

\item \textbf{Multisig key management.} Production vaults would use multisig (e.g.,
  2-of-3 cold keys). The interaction between threshold signing schemes and covenant
  enforcement is not addressed.

\item \textbf{Dust limit changes.} The per-UTXO recovery-cost scaling analysis assumes current
  dust limits (${\sim}546$~sats). Changes to dust relay policy would shift the
  economic floor at which splits become unrecoverable.
\end{itemize}

These threats are orthogonal to the covenant-specific properties measured in this
work. They represent important deployment considerations but do not affect the
relative ranking of covenant designs.

% ======================================================================
\section{Mitigations} \label{se:mitigations}

\subsection{TRUC/v3 Transactions}

TRUC (Topologically Restricted Until Confirmation), also known as v3 transactions,
limits descendant chains to a single transaction. TRUC was merged into Bitcoin Core
28.0~\cite{CoreRelay} and is available on mainnet as an opt-in relay policy.
If CTV vault transactions adopt TRUC, the 25-descendant fee pinning vector (TM1)
is eliminated at the relay policy level. This would remove CTV's primary
vulnerability, potentially shifting its
security profile: without fee pinning, CTV's worst-case attack becomes address reuse
(TM7), which is a user error preventable by wallet software rather than an adversarial
exploit. The TM1 row of the master threat matrix (and hence the CTV
branch of Proposition~\ref{thm:fee_inversion}'s counterexample
enumeration) would need re-evaluation under TRUC, as CTV would no
longer have a fee-invariant worst case under~TM1.

\subsection{Batched Recovery}

Class \classscale{} (TM4) is materially mitigated by batching
recovery transactions. Instead of recovering each split UTXO
individually, the watchtower combines multiple UTXOs into a single
recovery transaction, amortising fixed overhead across inputs. We
measure batched recovery on regtest for both CCV and OP\_VAULT
(Section~\ref{ss:batched_defender}); the fitted linear scaling is
$54 + 68 N$\vB{} for CCV and $154 + 92 N$\vB{} for OP\_VAULT. The
dust-threshold fee rate~$r^\dagger$ rises by $1.8\times$~(CCV) and
$2.7\times$~(OP\_VAULT), and the CCV/OP\_VAULT ordering under TM4
reverses relative to the single-input case.

\subsection{Poelstra Reset Model}

CAT+CSFS's cold key vulnerability (TM11) can be mitigated by replacing the bare
\texttt{OP\_CHECKSIG} recovery leaf with a recursive staging reset, as proposed by
Poelstra~\cite{Poe21}. Under this design, cold key compromise initiates a liveness
battle: the attacker resets the vault, the owner triggers a new withdrawal, the
attacker resets again. Each round costs both parties transaction fees. This converts
immediate theft into sustained denial of service, matching OP\_VAULT's security
profile under dual-key compromise. The tradeoff is a larger recovery script and
additional witness overhead.

% ======================================================================
\section{Open Problems} \label{se:open_problems}

\paragraph{Temporal verification.}
The Alloy models verify structural reachability but not timing properties. A TLA+
or UPPAAL model could capture the recovery race dynamics: given a CSV delay of
$d$~blocks, a block interval distribution, and network propagation delays, what is
the probability that the watchtower's recovery confirms before the attacker's
withdrawal? This would complement the structural verification with probabilistic
timing guarantees.

\paragraph{Cryptographic verification.}
The Alloy models abstract signatures as key-possession checks. A Tamarin or ProVerif
model could verify the cryptographic binding properties: that the CAT+CSFS dual
verification is sound under the Schnorr signature scheme, that CTV template hashes
are collision-resistant under SHA-256, and that Simplicity's jet-based verification
is equivalent to the combinator semantics. O'Connor's Coq formalization of
Simplicity~\cite{OC17} provides a foundation for this direction.

\paragraph{Cluster mempool impact.}
The ongoing Bitcoin Core cluster mempool redesign replaces per-transaction
ancestor/descendant limits with cluster-based evaluation~\cite{WD24}. This could
fundamentally change the fee pinning calculus: the 25-descendant limit exploited in
TM1 may be replaced by cluster size limits with different eviction logic. Measuring
the fee pinning attack under cluster mempool rules would determine whether CTV's
primary vulnerability survives the policy change.

\paragraph{Replacement cycling.}
Riard~\cite{Ria23} demonstrated transaction replacement cycling attacks against
Lightning Network HTLCs. The same class of attack may apply to covenant vaults:
an attacker could cycle the watchtower's recovery transaction out of the mempool
by repeatedly replacing it with conflicting transactions. This interaction between
covenant security and mempool replacement policy is unexplored.

\paragraph{Multi-vault composition.}
Real-world deployments would use multiple vaults simultaneously. The interaction
between concurrent vault operations---fee wallet contention in OP\_VAULT, UTXO set
management across vaults, and cross-vault batching opportunities---introduces
composition effects not captured by single-vault analysis.
